\subsection{VDA16 affine element-wise checkpoint}
\label{sec:vda16_affine_ew}

We evaluated one VDA16 affine element-wise checkpoint ($d_\mathrm{mem}=128$,
memory decay $=1.0$, seed 0) at terminal iteration 19,999 using the same
deterministic psychometric, recurrent-decoder, and attention-key clamp protocols
as the VDA analysis series (Fig.~\ref{fig:vda16_summary}). At displayed
validity 0.75 and the sampled orientation change nearest $18^\circ$, the
qualifying response rates were 0.893 for valid and
0.480 for forced-invalid changes
($\Delta=+0.413$). These are finite-trial estimates
from one checkpoint, not training-seed estimates.

\begin{figure*}[t]
  \centering
  \includegraphics[width=\textwidth]{figures/manuscript/vda16_manuscript_summary.pdf}
  \caption{VDA16 affine element-wise analysis. (A) Matched-bank response rates.
  (B) Foldwise train-only PCA-128 decoding from recurrent state.
  (C--D) Sensitivity and criterion under additive cued-location attention-key
  logit clamps. The checkpoint is competence-gated because its curriculum
  remained at $\theta=47.0^\circ$. VDA16
  d256, decay, and replicated-seed cells are unavailable.}
  \label{fig:vda16_summary}
\end{figure*}

\begin{figure*}[t]
  \centering
  \includegraphics[width=\textwidth]{figures/manuscript/vda_method_comparison_with_vda16.pdf}
  \caption{Descriptive d128 routing-family matrix. The open, unconnected
  VDA16 point is competence-gated and is not part of a fitted or implied
  set-size trend. VDA2 is absent because no compatible modern matched-width
  shard exists. Historical and VDA16 cells use separately frozen
  executable/runtime lineages.}
  \label{fig:vda16_method_context}
\end{figure*}

\begin{figure*}[t]
  \centering
  \includegraphics[width=\textwidth]{figures/clamp/vda16_clamp_achieved_attention.pdf}
  \caption{Achieved cued-location attention mass under the VDA16 logit-clamp
  series. Self-attention key mass is averaged over query tokens.
  Trial-level error bars describe evaluation trials
  from one checkpoint, not training-seed uncertainty.}
  \label{fig:vda16_achieved_attention}
\end{figure*}
